\chapter{Коллективные координаты и статистика бозонного дефекта}

Нейтральное поле $Q$ допускает конечный проекторный дефект, но его
частичный тип определяется не только статической топологией. Необходимо
установить, какие деформации являются нормируемыми нулевыми модами и
допускает ли их пространство нетривиальное квантование
Финкельштейна--Рубинштейна.

\section{Формальная ориентационная орбита}

До выбора ежа симметрия имеет вид
\begin{equation}
 SO(3)_{\rm space}\times SO(3)_{\rm fam}.
\end{equation}
Профиль
\begin{equation}
 Q_0(x)=q(r)\left(\hat x\hat x^T-\frac13I_3\right)
\end{equation}
инвариантен относительно диагонального $SO(3)$. Поэтому формальное
пространство относительных ориентаций равно $SO(3)$, а его фундаментальная
группа есть $\mathbb Z_2$.

Однако размерность орбиты ещё не доказывает существование локальной
коллективной координаты. Её норма должна быть конечной в кинетической
метрике поля.

\section{Кубическая расходимость момента инерции}

Для внутреннего вращения с генератором $\Omega$ вариация равна
\begin{equation}
 \delta_\Omega Q=[\Omega,Q].
\end{equation}
Для единичного вращения вокруг третьей оси угловой интеграл имеет точное
значение
\begin{equation}
 \int_{S^2}\Tr([\Omega,P]^T[\Omega,P])\,d\Omega_{S^2}
 =\frac{16\pi}{3}.
\end{equation}
При профиле
\begin{equation}
 q(r)=\Delta\frac{r^2}{r^2+R^2}
\end{equation}
момент инерции до радиуса $L$ равен
\begin{equation}
 \mathcal I(L)=\frac{16\pi}{3}\int_0^Lr^2q(r)^2dr
 =\frac{16\pi}{9}\Delta^2L^3+O(L).
 \label{eq:v6-bosonic-orientation-inertia}
\end{equation}

Численная аппроксимация при $L=40,80,160$ дала степень
$3.002412\ldots$. Предсказанный коэффициент $L^3$ равен
$4.2103367\ldots$, измеренный --- $4.2093620\ldots$.

\begin{theorem}[Отсутствие локального внутреннего ротора]
Относительное семейное вращение меняет упорядоченный вакуум во всём бесконечном
объёме и имеет кубически расходящийся момент инерции. Поэтому формальная
орбита $SO(3)$ не является нормируемым пространством коллективных
координат одиночного дефекта.
\end{theorem}

Одновременное пространственное и семейное вращение даёт нулевую вариацию:
это стабилизатор конфигурации, а не новый модуль. Машинный остаток данного
тождеcтва равен нулю.

\section{Оставшиеся локальные моды}

Три переноса центра являются точными статическими модулями вследствие
однородности пространства. Их окончательная кинетическая норма требует
лоренц-ковариантного временного продолжения составной связности, но
конечность статической энергии оставляет этот маршрут открытым.

Масштаб не является нулевой модой. Для положительных вкладов
\begin{equation}
 E(R)=c_D R+\frac{c_F}{R}+c_VR^3
\end{equation}
вторая производная в стационарной точке положительна:
\begin{equation}
 E''(R_*)=\frac{2c_F}{R_*^3}+6c_VR_*>0.
\end{equation}
При единичных коэффициентах пробного профиля получены
$R_*=0.8301754\ldots$ и кривизна $204.4017\ldots$. Это демонстрирует
массивную пульсационную моду, но не задаёт её физическую частоту, поскольку
общая временная нормировка ещё не выведена.

Таким образом, в текущем приближении локализованное пространство модулей
сводится к $\mathbb R^3$ переносов и имеет тривиальную фундаментальную
группу.

\section{Почему нет дионной фазы}

Для обычного единичного BPS-монополя пространство модулей имеет вид
$\mathbb R^3\times S^1$; окружность возникает из глобальной фазы
остаточной калибровочной группы и после квантования несёт электрический заряд.
См. \path{arXiv:hep-th/9603086} и
\path{arXiv:hep-th/9506052}.

В проекте составная связность $A_Q$ не содержит стабилизаторное
$O(2)$-направление. Поэтому аналогичной нормируемой $S^1$-координаты нет,
и дионоподобный электрический заряд не возникает.

\section{Статистика Финкельштейна--Рубинштейна}

Формальная орбита $SO(3)$ действительно допускает два одномерных
представления $\pi_1(SO(3))=\mathbb Z_2$: тривиальное бозонное и знаковое
фермионное. Общая связь вращения и обмена солитонов исследована в
\path{Commun.Math.Phys.115(1988)421-434}; исходный механизм --- DOI
\path{10.1063/1.1664510}.

Но не нормируемая ориентация не входит в локальное пространство модулей.
Полная Real WZW/Pfaffian-фаза проекта равна $+1$ и не отождествлена с
голономией $2\pi$-вращения. Поэтому нет выведенного нетривиального
ограничения Финкельштейна--Рубинштейна.

\begin{proposition}[Минимальное коллективное квантование]
В текущем родителе основное состояние одиночного нейтрального дефекта
квантуется как скалярный бозон. Полуцелый спин не следует из его
нормируемых коллективных координат. Это не запрещает будущую новую
топологическую линию над полным пространством конфигураций, но такая линия
сейчас отсутствует.
\end{proposition}

Следующая проверка должна определить, фиксирует ли единый родитель
коэффициенты $c_D,c_F,c_V$, абсолютные массу и радиус, а также хотя бы один
ненулевой портал к наблюдаемому сектору. Без этого нейтральный солитон
остаётся математически локализованным, но феноменологически ненормированным.

\begin{protocol}
\begin{enumerate}
 \item формальная относительная орбита: \textbf{$SO(3)$};
 \item диагональное вращение: \textbf{стабилизатор};
 \item момент инерции внутреннего вращения: \textbf{$\sim L^3$};
 \item нормируемый внутренний ротор: \textbf{нет};
 \item локальные трансляционные моды: \textbf{три};
 \item масштабная нулевая мода: \textbf{нет};
 \item массивная пульсационная мода: \textbf{условно есть};
 \item дионоподобная $S^1$-фаза: \textbf{нет};
 \item нетривиальный FR-знак: \textbf{не выведен};
 \item минимальный спин основного состояния: \textbf{нуль};
 \item текущая статистика: \textbf{бозонная};
 \item абсолютные масса и частота: \textbf{не получены};
 \item следующая проверка: \textbf{родительская масса и портал}.
\end{enumerate}
\end{protocol}

Машинный сертификат сохранён в
\path{s2t_v6_bosonic_defect_collective_quantization_gate_results.json}.